\mxunit{automatismes-ch05}{Statistiques descriptives à une variable}

\begin{fluencysheet}{jeu.~17~déc.~2026}{Moyenne et linéarité de la moyenne}
\mxfq{Échauffement · Flash}{STA}{Moyenne de $12,\,15,\,9,\,14$\,?}{$\dfrac{50}{4}=\num{12,5}$.}{}
\mxfq{Cœur · Méthode}{STA}{La série $4,\,6,\,8,\,10,\,12$ a pour moyenne $8$. Quelle est la moyenne si l'on ajoute $3$ à chaque valeur\,? si on multiplie chaque valeur par $2$\,?}{Linéarité~: $8+3=11$\,; $2\times8=16$.}{Pas besoin de refaire la somme.}
\mxfq{Cœur + · Vrai/faux}{STA}{Vrai ou faux\,? Justifier. La moyenne de deux nombres $a<b$ est toujours strictement comprise entre $a$ et $b$.}{Vrai. $a<\dfrac{a+b}{2}<b$ car $b-a>0$.}{}
\mxfq{Retour · Lecture}{LEC}{Dans un diagramme circulaire de $80$ personnes, un secteur représente $\pc{25}$. Quel est son effectif\,? Quel est son angle\,?}{$20$ personnes\,; $90\degre$.}{}
\mxfq{Fil rouge · Flash}{DIV}{Calculer $\pgcd\pt{84;126}$ et en déduire la forme irréductible de $\dfrac{84}{126}$.}{$\pgcd=42$\,; $\dfrac{84}{126}=\dfrac23$.}{}
\end{fluencysheet}

\begin{fluencysheet}{lun.~4~janv.~2027}{Médiane et quartiles}
\mxfq{Échauffement · Flash}{STA}{Ranger $7,\,2,\,9,\,4,\,5,\,11,\,3$ puis donner la médiane.}{$2\,;\,3\,;\,4\,;\,5\,;\,7\,;\,9\,;\,11$\,; médiane $5$ (rang $4$).}{}
\mxfq{Cœur · Méthode}{STA}{Série $3\,;\,5\,;\,6\,;\,8\,;\,9\,;\,12\,;\,14\,;\,15$. Donner $Q_1$, la médiane, $Q_3$ et l'écart interquartile.}{$n=8$~: $Q_1$ rang $\lceil8/4\rceil=2$~: $5$\,; médiane $\dfrac{8+9}{2}=\num{8,5}$\,; $Q_3$ rang $6$~: $12$\,; $Q_3-Q_1=7$.}{Convention~: plus petite valeur telle qu'au moins $\pc{25}$ (resp. $\pc{75}$) des données lui soient inférieures ou égales.}
\mxfq{Cœur + · Vrai/faux}{STA}{Vrai ou faux\,? Justifier. Au moins la moitié des valeurs d'une série sont supérieures ou égales à sa médiane.}{Vrai. Par définition de la médiane (au moins $\pc{50}$ des valeurs de chaque côté, bornes comprises).}{}
\mxfq{Retour · Flash}{CMP}{Ranger $\num{0,6}$, $\dfrac58$ et $\pc{61}$ dans l'ordre croissant.}{$\num{0,6}<\pc{61}=\num{0,61}<\dfrac58=\num{0,625}$.}{}
\mxfq{Fil rouge · Graphique → formule}{REG}{La droite ci-dessous représente une fonction affine $f$. Déterminer $f(x)$.\par\smallskip\begin{tikzpicture}[scale=0.45,>=stealth]
\draw[very thin,gray!45] (-1,-1) grid (5,3);
\draw[->] (-1,0)--(5.3,0); \draw[->] (0,-1)--(0,3.3);
\foreach \i in {-1,1,2,3,4,5}\draw (\i,.1)--(\i,-.1) node[below,font=\tiny]{\i};
\foreach \j in {-1,1,2,3}\draw (.1,\j)--(-.1,\j) node[left,font=\tiny]{\j};
\draw[thick,mxink,domain=-1:5,samples=60,smooth] plot (\x,{-0.5*\x+2});
\fill[mxteacher] (0,2) circle (3pt);
\fill[mxteacher] (2,1) circle (3pt);

\end{tikzpicture}\par }{Elle passe par $\pt{0;2}$ et $\pt{2;1}$~: pente $\dfrac{1-2}{2-0}=-\dfrac12$, $f(x)=-\dfrac12x+2$.}{}
\end{fluencysheet}

\begin{fluencysheet}{mar.~5~janv.~2027}{Écart type\,; effet de l'ajout ou de la suppression d'une valeur}
\mxfq{Échauffement · Flash}{CAL}{Calculer $(3-5)^2+(7-5)^2$.}{$4+4=8$.}{}
\mxfq{Cœur · Méthode}{STA}{Série $2\,;\,4\,;\,6$~: moyenne, variance et écart type (valeur approchée à $10^{-2}$).}{Moyenne $4$\,; variance $\dfrac{4+0+4}{3}=\dfrac83$\,; écart type $\sqrt{\dfrac83}\approx\num{1,63}$.}{}
\mxfq{Cœur + · Vrai/faux}{STA}{Vrai ou faux\,? Justifier. Ajouter $5$ à chaque valeur d'une série ne change pas son écart type.}{Vrai. Toutes les valeurs et la moyenne sont décalées de $5$~: les écarts à la moyenne sont inchangés.}{}
\mxfq{Retour · Programme de construction}{CON}{Soit $ABC$ un triangle, $I$ le milieu de $[AB]$ et $D$ le symétrique de $C$ par rapport à $I$. Quelle est la nature de $ACBD$\,?}{Un parallélogramme~: ses diagonales $[AB]$ et $[CD]$ ont même milieu $I$.}{}
\mxfq{Fil rouge · Flash}{GEO}{$(MN)\parallel(BC)$, $M\in[AB]$, $N\in[AC]$, $AM=3$, $AB=5$, $AC=10$. Calculer $AN$.}{Thalès~: $\dfrac{AN}{AC}=\dfrac{AM}{AB}$, donc $AN=10\times\dfrac35=6$.}{}
\end{fluencysheet}

\begin{fluencysheet}{jeu.~7~janv.~2027}{Comparaison de séries}
\mxfq{Échauffement · Flash}{STA}{Étendue de $3,\,8,\,5,\,12$\,?}{$12-3=9$.}{}
\mxfq{Cœur · Interpréter}{STA}{Deux séries $A$ et $B$ ont même moyenne $12$\,; $\sigma_A=\num{1,5}$ et $\sigma_B=4$. Laquelle est la plus régulière\,?}{$A$~: les valeurs sont plus resserrées autour de la moyenne ($\sigma$ plus petit).}{}
\mxfq{Cœur + · Interpréter}{STA}{$A$~: médiane $11$, écart interquartile $2$. $B$~: médiane $13$, écart interquartile $6$. Comparer les deux séries.}{$B$ est centrée plus haut (médiane $13>11$) mais plus dispersée (écart interquartile $6>2$).}{}
\mxfq{Retour · Flash}{PRB}{On lance deux pièces équilibrées. Probabilité d'obtenir deux faces\,?}{$4$ issues équiprobables~: $\dfrac14$.}{}
\mxfq{Fil rouge · Méthode}{ALG}{Programme~: choisir un nombre\,; l'élever au carré\,; ajouter $1$. Résultat pour $-3$\,? Peut-on obtenir $0$\,? Justifier.}{Pour $-3$~: $(-3)^2+1=10$. Non~: $x^2\geqslant0$ donc $x^2+1\geqslant1>0$.}{}
\end{fluencysheet}

\begin{fluencysheet}{lun.~11~janv.~2027}{Données regroupées en classes~: histogramme et fréquences cumulées}
\mxfq{Échauffement · Flash}{STA}{Donner l'amplitude et le centre de la classe $\intfo{20}{35}$.}{Amplitude $15$\,; centre $\dfrac{20+35}{2}=\num{27,5}$.}{}
\mxfq{Cœur · Tableau}{LEC}{Classes $\intfo{0}{10}$, $\intfo{10}{20}$, $\intfo{20}{30}$ d'effectifs $4,\,6,\,10$. Donner les fréquences puis les fréquences cumulées croissantes.}{Total $20$\,; fréquences $\pc{20}$, $\pc{30}$, $\pc{50}$\,; cumulées $\pc{20}$, $\pc{50}$, $\pc{100}$.}{}
\mxfq{Cœur + · Flash}{LEC}{Dans un histogramme à classes de même amplitude, deux rectangles ont pour hauteurs $3$ cm et $6$ cm. Quel est le rapport de leurs effectifs\,?}{Les hauteurs sont proportionnelles aux effectifs~: $1$ pour $2$.}{}
\mxfq{Retour · Flash}{NUM}{Écrire $\num{0,00032}$ en notation scientifique.}{$\num{3,2}\times10^{-4}$.}{}
\mxfq{Fil rouge · Flash}{EQ}{Résoudre dans $\R$~: $2(x+3)=5x-9$.}{$x=5$}{Piège~: oublier de distribuer le $2$ sur le $3$.}
\end{fluencysheet}

\begin{fluencysheet}{mar.~12~janv.~2027}{Moyenne pondérée, classe médiane\,; tableur et Python}
\mxfq{Échauffement · Flash}{CAL}{Calculer $\dfrac{2\times8+3\times12}{5}$.}{$\dfrac{52}{5}=\num{10,4}$.}{}
\mxfq{Cœur · Méthode}{STA}{Classes $\intfo{0}{10}$, $\intfo{10}{20}$, $\intfo{20}{30}$ d'effectifs $4\,;\,7\,;\,9$. Estimer la moyenne avec les centres de classes\,; donner la classe médiane.}{Moyenne estimée $\dfrac{4\times5+7\times15+9\times25}{20}=\num{17,5}$\,; effectifs cumulés $4\,;\,11\,;\,20$~: les $10^{\text e}$ et $11^{\text e}$ valeurs sont dans $\intfo{10}{20}$~: classe médiane.}{}
\mxfq{Cœur + · Lecture de code}{ALG}{Que calcule ce programme et quelle valeur obtient-on\,?\mxpseudo{notes = [12, 15, 9, 14]\\ s = 0\\ for n in notes:\\ \mxtab s = s + n\\ print(s / len(notes))}}{La moyenne des notes~: $\dfrac{50}{4}=\num{12,5}$.}{}
\mxfq{Retour · Flash}{PRO}{Un article à $\num{45000}$ Ar est soldé à $-\pc{30}$. Quel est le prix soldé\,?}{$\num{45000}\times\num{0,70}=\num{31500}$ Ar.}{}
\mxfq{Fil rouge · Prouver}{LIT}{Programme~: choisir $n$\,; ajouter $1$\,; élever au carré\,; soustraire $n^2$. Montrer que le résultat est toujours un nombre impair.}{$(n+1)^2-n^2=2n+1$~: impair pour tout entier $n$.}{}
\end{fluencysheet}
